\scanpage{45}
\mspage{43}
% archive sheet number: 44
au dessus duquel $\inserted{\unsure{E_\alpha|\mathscr E_\alpha}=}R^if_{\alpha *}(\unsure{\mathbb Z_{\ell\,X_\alpha}|Z_{\ell\,X_\alpha}})$ est constant tordu. En effet, le fait que \deleted{$\unsure{E}$} $R^if_{K*}(\mathbb Z_{\ell\,X_K})$ s'en déduise par le chgt de base $\eta\longrightarrow\eta_\alpha\longrightarrow S_\alpha$ signifie que la formation de $R^if_{\alpha *}(\mathbb Z_{\ell\,X_\alpha})$ commute avec les chgts de base $\eta_\alpha-S_\alpha$ et $\eta-\eta_\alpha$. Pour le premier c'est trivial, pour le deuxième c'est une conséquence du th. de chgt de base \deleted{lisse} par morphisme lisse (SGA 4 \uline{XVI} 1.6). \deleted{D'autre part, l'exis}
D'autre part, l'existence d'un $\unsure{V_\alpha|U_\alpha}$ comme dessus est \unsure{encore prouvée|encore garantie}, \unsure{disposant de|si on dispose de} la résolution des singularités \inserted{(sous la forme de Hironaka)} pour les schémas de type fini et de dimension $\leq\dim X_{K\deleted{\unsure{\alpha}}}$ sur $K_\alpha$ (SGA 5 \uline{II} 5.1). \deleted{\unsure{On le voit ainsi}} \inserted{Cette hypothèse est vérifiée} en particulier, par Hironaka, si $K$ est de car. nulle, et \deleted{\unsure{justifie le th. de monodromie lorsque les termes dans}} \inserted{par Abhyankar si $\dim X_K\leq2$.}

\deleted{\uline{Corollaire} 8.9. Soit $V$ comme dans 8.1, \inserted{$X$ de type fini sur $X_K$.} \deleted{\unsure{sur $K=k(V)$}}, \deleted{\unsure{de caractéristique}}, \inserted{\unsure{hens.} i.e. $\mathbb Z$,} \uline{caractéristique} \inserted{\uline{nulle}, ou $\dim X_K\leq2$}, \deleted{soit $X_K$ de t.f. sur $K$}, \deleted{\unsure{Supp.}} \deleted{\unsure{pour}} et \unsure{supposons} les $V_\alpha$ de valuation discrète (cas de 8.2. a) p.ex) ou l'action de \deleted{\unsure{resp.}} $I$ dans $H^i(X_{\overline K},\mathbb Z_\ell)$ modérée (cas de 8.2 b) par exemple). Alors l'\deleted{\unsure{représentation}} \inserted{action} de $I$ sur $H^i(X_{\overline K},\mathbb Z_\ell)$ est quasi-unipotente.}
\margin{\deleted{\uline{Supposons $K$ de}}}

\uline{Corollaire} 8.9. Les conclusions de 8.2 restent vraies en supposant seulement $X$ de t.f. sur $K$ (pas néc. propre), pourvu que dans le cas a), on suppose la caractéristique de $K$ égale à $0$, \uline{ou} $\dim X_K\leq2$.
